\section{Field Names}\label{sec:Fields}
I distinguish generally two types of fields\index{field} that will have different naming rules.

\paragraph{Attributes}\index{field!attribute} are mutable or immutable fields\index{field}, that get their values during runtime; that they may get assigned a default value at compile time that may never change during runtime is irrelevant.

\paragraph{Constants}\index{field!constant} are usually static final fields\index{field} that get their values always during compile time. They should have an immutable type, like a primitive or \lstinline|java.lang.String|\index{java.lang!String}.

%------------------------------------------------------------------------------

\subsection{Names for Attributes}\label{sec:AttributeNames}
The identifier for an attribute\index{field!attribute} has to be prefixed with ``\verb#m_#'', usually followed by the property\index{property} name (see \tqfullvref{sec:NamesOfProperties}) with the first letter capitalised:
\begin{lstlisting}!local
private String m_AString;
private boolean m_IsValid;https://docs.oracle.com/en/java/javase/14/language/records.html
private String m_HTMLHeader;
private MessageProvider m_MessageProvider;
\end{lstlisting}
This ensures that the names of local variables and formal parameters of methods\index{method} or constructors\index{constructor} will be always distinct from those of attributes\index{field!attribute}, and therefore neither local variables nor parameters can collide with or hide an attribute\index{field!attribute}. Another consequence is that it is not necessary to use \lstinline|this.| when accessing an attribute\index{field!attribute}.

Aside this, anything else that was said about the naming of local variables in chapter \tqref{sec:NamesForLocalVariables} is also valid for the naming of attributes\index{field!attribute}.

Of course not all attributes\index{field!attribute} of a class\index{class} are properties\index{property}. They could hold an internal status as well or may serve other purposes. But the naming rules remain the same:  the prefix ``\verb#m_#'' is followed by a name describing the purpose of the attribute\index{field!attribute}, typed in \textit{CamelCase} with the first letter capitalised.

Eclipse\index{Eclipse} knows a configuration setting for the attribute\index{field!attribute} https://docs.oracle.com/en/java/javase/14/language/records.htmlname prefix (\verb#Window|Preferences|Java|Code Style#). Setting the prefix list there to “\verb#m_#” makes sure that the generation of getters\index{method!getter method} and setters\index{method!setter method} will work as expected.

For IntelliJ~IDEA\index{IntelliJ IDEA}, a similar setting can be found at \verb#File|Settings#, where you select under \verb#Editor|Code Style|Java# the tab \verb#Code Generation#.

%------------------------------------------------------------------------------

\subsection{Names for Constants}\label{sec:Constants}
The \Java programming language defines a keyword \lstinline|const|\index{const} \autocite{ORACLE_DOC_LANGUAGE_SPECIFICATION:Keywords}, but that is just a reserved word without any function. There is also no syntactical concept for constant\index{constant} values in Java.

This means that a constant\index{constant} in Java exists only by convention. 

Constants\index{constant} in Java are fields\index{field!constant} of primitive or immutable types\index{immutable type} that are declared as public, static and final.

Our definition for a constant\index{field!constant} is even more restrictive: that public static final field\index{field} has to be initialised with a compile-time constant\index{constant} value. This basically limits constants\index{field!constant} to \bdq{magic numbers}\index{magic number} and texts.\footnote{This limitation is imposed by the documentation rules, refer to \tqfullvref{sec:FieldComment}.}

According to the \ccjpl, the name of a constant\index{field!constant} ``should be all uppercase with words separated by underscores (`\verb#_#')'' \autocite{SUN_CODE_CONVENTIONS:NamingConventions}.

\keeplisting{5}
Some samples for valid constants\index{field!constant}
\begin{lstlisting}
public static final double PI = 3.1415;
public static final int ANSWER_TO_ALL_QUESTIONS = 42;
public static final String ISO_DATE_FORMAT = "yyyyMMdd'T'hhmmss";
\end{lstlisting}

Real constants\index{field!constant} go to the section of the class headlined with \bdq{Constants} (refer to chapter \tqfullvref{sec:StructuringComments}).

\keeplisting{4}
The field \lstinline|PATTERN| below is not a constant\index{field!constant} in the sense of our definition, because it is not initialised with a compile-time constants:
\index{java.util.regex!Pattern!compile()}
\begin{lstlisting}
// NOT A CONSTANT
public static final Pattern PATTERN = Pattern.compile( ".*" );
\end{lstlisting}

\keeplisting{5}
It should not get named as a constant\index{field!constant}, and it should go to the \bdq{Static Initialisations} section of the class\index{class} (see \tqref{sec:StructuringComments}). In most cases, it shouldn't be public, either.
\begin{lstlisting}
    /*------------------------*\
====** Static Initialisations **=====================================
    \*------------------------*/
private static final Pattern m_Pattern = Pattern.compile( ".*" );
\end{lstlisting}

You will encounter this kind of code pattern (sorry for the duplication of terms!) quite often, and the best way to implement it would look like this (for an instance of \lstinline|java.util.regex.Pattern| \autocite{ORACLE_DOC_PATTERN_CLASS}):

\keeplisting{21}
\index{java.lang!String}\index{java.util.regex!Pattern}\index{java.util.regex!Pattern!compile()}\index{java.util.regex!PatternSyntaxException}\index{java.lang!ExceptionInInitializerError}
\begin{lstlisting}
    /*-----------*\
====** Constants **==================================================
    \*-----------*/
public static final String PATTERN_SOURCE = ".*";

    /*------------------------*\
====** Static Initialisations **=====================================
    \*------------------------*/
private static final Pattern m_Pattern;

static
{
    try
    {
        m_Pattern = Pattern.compile( PATTERN_SOURCE );
    }
    catch( final PatternSyntaxException e )
    {
        throw new ExceptionInInitializerError( e );
    }    
}
\end{lstlisting}

In case you have a constant\index{constant} value (a magic number\index{magic number}, a message string or alike) that should not be public for some reason, you can name that either as a regular field\index{field!attribute} (an attribute with the ``\verb#m_#'' prefix) or as a regular constant\index{field!constant} as described above. For a local constant\index{local constant} (a local variable as an alias for a magic number\index{magic number} or a message string~…), use the naming pattern for constants\index{field!constant}. 

Sometimes you may have classes, groups, families or categories of constants\index{constant} (e.g. the tags in an XML\index{XML} document, the column names of a database table, or message texts). It has proved that it supports the readability of the code if this is reflected in the names for those constants\index{field!constant}. For the column names of the table \verb#PERSON#, the constants\index{field!constant} may be defined as shown in this example:
\keeplisting{4}
\index{java.lang!String}
\begin{lstlisting}
public final static String PERSON_COL_FirstName = "first_name";
public final static String PERSON_COL_LastName = "last_name";
public final static String PERSON_COL_MaidenName = "maiden_name";
public final static String PERSON_COL_Title = "title";
…
\end{lstlisting}
The \textit{camelCased} suffix is the \bsq{real} name, while the part that follows the stronger rules for a constant\index{field!constant} name is just indicating the \bdq{constant category}, \verb#PERSON_COL# in this case.

%------------------------------------------------------------------------------

\subsection{Names for Record Components}\label{sec:NamesOfRecordComponents}
Technically, the record\index{record} components are attributes\index{field!attribute} and should be named as described in \tqref{sec:AttributeNames}:
\begin{lstlisting}
// NOT RECOMMENDED!
public record MyRecord( int m_Number, String m_Text ) {};
\end{lstlisting}
But this would mean that the accessor methods\index{method!accessor method} would be named as \lstinline|m_Number()| and \lstinline|m_Text()|, and that is confusing when used.

Additionally, the record\index{record} declaration resembles a constructor\index{constructor}.

Consequently, we use the naming rules for parameters from \tqfullref{sec:NamesForFormalParameters} for the record components.
\begin{lstlisting}
// RECOMMENDED!
public record MyRecord( int number, String text ) {};
\end{lstlisting}

%------------------------------------------------------------------------------

\subsection{Summary}

\subsubsection{Principles}
\begin{enumerate}[label=P\arabic*.]
	\item{Attributes\index{field!attribute} and constants\index{field!constant} are both fields\index{field}, but with different naming rules.
	
	Attributes\index{field!attribute} basically are fields\index{field} that get their value during runtime, while constants are public static final fields\index{field}.
	
	Constants\index{field!constant} are fields that are \textit{public}, \textit{static} and \textit{final}, and are initialised with a compile-time constant\index{constant} value.}
	\item{The name for an attribute\index{field!attribute} is prefixed with ``\verb#m_#'', followed by the name in \textit{CamelCase}, with an uppercase first letter.}

	\item{The names of constants\index{field!constant} are all uppercase, with underscores (\verb#_#) as separators between words.
	
	To group constants\index{field!constant} it is possible to name them using a regular constant name (all uppercase~…) followed by an underscore as prefix, and the name in CamelCase.}
	
	\item{Usually use the singular, but for collections and arrays, use the plural form.}
	
	\item{Record\index{record} components are named like parameters, not like attributes\index{field!attribute}.}
\end{enumerate}

%==============================================================================
% End of file!!
%==============================================================================
